% ══════════════════════════════════════════════════════════════════════
% StenoWav — Technical Reference & Mathematical Formalisation
% ══════════════════════════════════════════════════════════════════════
\documentclass[11pt,a4paper]{article}

\usepackage[utf8]{inputenc}
\usepackage[T1]{fontenc}
\usepackage{amsmath,amssymb,amsthm}
\usepackage{mathtools}
\usepackage{algorithm}
\usepackage{algpseudocode}
\usepackage{hyperref}
\usepackage{geometry}
\usepackage{booktabs}
\usepackage{enumitem}
\usepackage{tikz}
\usepackage{xcolor}
\usepackage{listings}
\usepackage{fancyhdr}
\usepackage{tcolorbox}

\geometry{margin=2.5cm, headheight=14.5pt}
\hypersetup{colorlinks=true, linkcolor=blue!60!black, citecolor=blue!60!black, urlcolor=blue!60!black}

\pagestyle{fancy}
\fancyhf{}
\rhead{\textit{StenoWav Technical Reference}}
\lhead{\textit{\nouppercase{\leftmark}}}
\rfoot{\thepage}

% Theorem environments
\newtheorem{theorem}{Theorem}[section]
\newtheorem{lemma}[theorem]{Lemma}
\newtheorem{proposition}[theorem]{Proposition}
\newtheorem{corollary}[theorem]{Corollary}
\newtheorem{definition}[theorem]{Definition}
\newtheorem{remark}[theorem]{Remark}

% Custom commands
\newcommand{\R}{\mathbb{R}}
\newcommand{\Z}{\mathbb{Z}}
\newcommand{\N}{\mathbb{N}}
\newcommand{\F}{\mathcal{F}}
\newcommand{\xor}{\oplus}
\newcommand{\floor}[1]{\left\lfloor #1 \right\rfloor}
\newcommand{\ceil}[1]{\left\lceil #1 \right\rceil}
\newcommand{\abs}[1]{\left| #1 \right|}
\newcommand{\norm}[1]{\left\| #1 \right\|}
\DeclareMathOperator{\round}{round}
\DeclareMathOperator{\argmin}{argmin}
\DeclareMathOperator{\rms}{RMS}
\DeclareMathOperator{\snr}{SNR}

\lstset{
    basicstyle=\ttfamily\small,
    breaklines=true,
    frame=single,
    framesep=3pt,
    backgroundcolor=\color{gray!8},
}

\title{
    \textbf{StenoWav} \\[0.5em]
    \Large Technical Reference \& Mathematical Formalisation \\[0.3em]
    \normalsize Version 0.1.0
}
\author{StenoWav Project}
\date{March 2026}

\begin{document}
\maketitle
\tableofcontents
\newpage

% ──────────────────────────────────────────────────────────────────
\section{Introduction}
% ──────────────────────────────────────────────────────────────────

StenoWav is an audio steganography framework that embeds arbitrary binary data
inside 16-bit PCM WAV audio signals.  The core embedding primitive is
\emph{Quantisation Index Modulation} (QIM), applied through a packet-based
architecture with per-packet cipher obfuscation.

The framework supports two embedding domains:
\begin{enumerate}
    \item \textbf{Time domain} --- direct modification of PCM samples.
    \item \textbf{Frequency domain} --- modification of DFT magnitude
          coefficients, with phase preservation and IDFT reconstruction.
\end{enumerate}

Both paths use RMS-adaptive segment analysis to restrict embedding to
signal regions with sufficient energy, and an optional dither pass to
raise the statistical noise floor.

This document provides the full mathematical formalisation of every
algorithm in the pipeline, along with proofs of key properties (distortion
bounds, extraction correctness, security guarantees).

\subsection{Notation}

\begin{center}
\begin{tabular}{@{}ll@{}}
\toprule
Symbol & Meaning \\
\midrule
$x \in \R$ & A single audio sample (normalised to $[-1, 1]$) \\
$\Delta \in \R^+$ & QIM step size (quantisation interval) \\
$q \in \Z$ & Quantisation index: $q = \round(x / \Delta)$ \\
$M \in \N$ & QIM modulus (number of lattice classes) \\
$s \in \{0, \ldots, M-1\}$ & Symbol to embed \\
$b \in \{0, 1\}$ & Single bit to embed (special case $M=2$) \\
$K \in \{0, \ldots, 2^{16}-1\}$ & 16-bit cipher key \\
$\sigma$ & 32-byte master seed \\
$p$ & Packet number ($p \in \N_0$) \\
$\F$ & Discrete Fourier Transform \\
\bottomrule
\end{tabular}
\end{center}


% ──────────────────────────────────────────────────────────────────
\section{Quantisation Index Modulation (QIM)}
\label{sec:qim}
% ──────────────────────────────────────────────────────────────────

\subsection{Standard QIM ($M = 2$)}

\begin{definition}[QIM Lattice]
Given step size $\Delta > 0$, define the quantisation lattice
$\Lambda = \{ k\Delta : k \in \Z \}$.
Partition $\Lambda$ into two sub-lattices indexed by parity:
\[
    \Lambda_b = \{ k\Delta : k \equiv b \pmod{2} \}, \quad b \in \{0,1\}.
\]
\end{definition}

\begin{definition}[QIM Embedding]
Given sample $x$ and bit $b$:
\begin{enumerate}
    \item Compute the quantisation index $q = \round(x / \Delta)$.
    \item If $q \bmod 2 = b$, set $q^* = q$ (no shift needed).
    \item If $q \bmod 2 \neq b$, pick the nearest index with correct parity:
    \[
        q^* = \begin{cases}
            q + 1 & \text{if } x \geq q\Delta, \\
            q - 1 & \text{otherwise.}
        \end{cases}
    \]
    \item Output $x' = q^* \Delta$.
\end{enumerate}
\end{definition}

\begin{definition}[QIM Extraction]
Given watermarked sample $x'$:
\[
    \hat{b} = \round(x' / \Delta) \bmod 2.
\]
\end{definition}

\begin{theorem}[Extraction correctness]
\label{thm:qim-correct}
For any $x \in \R$, $b \in \{0,1\}$, and $\Delta > 0$, let $x' = \mathrm{QIM}_\Delta(x, b)$.
Then $\round(x' / \Delta) \bmod 2 = b$.
\end{theorem}

\begin{proof}
By construction, $x' = q^*\Delta$ where $q^* \equiv b \pmod{2}$.
Since $\round(q^*\Delta / \Delta) = q^*$, the extracted parity is
$q^* \bmod 2 = b$.
\end{proof}

\begin{theorem}[Distortion bound]
\label{thm:distortion}
For any embedding $x' = \mathrm{QIM}_\Delta(x, b)$:
\[
    \abs{x' - x} \leq \Delta.
\]
\end{theorem}

\begin{proof}
In the worst case, $x$ lies exactly at $q\Delta$ with $q \bmod 2 \neq b$,
requiring a shift to $q^* = q \pm 1$, giving $\abs{x' - x} = \Delta$.
In all other cases the shift is smaller.
\end{proof}

\subsection{Multi-Symbol QIM ($M \geq 2$)}
\label{sec:mod-m-qim}

\begin{definition}[Mod-$M$ QIM]
Partition $\Lambda$ into $M$ sub-lattices:
$\Lambda_s = \{ k\Delta : k \equiv s \pmod{M} \}$ for $s \in \{0, \ldots, M-1\}$.

Embedding symbol $s$ into sample $x$:
\begin{enumerate}
    \item $q = \round(x / \Delta)$.
    \item $r = q \bmod M$ (using Euclidean remainder).
    \item If $r = s$, output $q\Delta$.
    \item Otherwise, compute the forward offset $d = (s - r) \bmod M$:
          \begin{align}
              q_{\uparrow} &= q + d, \\
              q_{\downarrow} &= q + d - M.
          \end{align}
    \item Pick the candidate closer to $x$:
          \[
              q^* = \argmin_{q' \in \{q_\uparrow, q_\downarrow\}} \abs{q'\Delta - x}.
          \]
    \item Output $x' = q^* \Delta$.
\end{enumerate}
\end{definition}

Extraction: $\hat{s} = \round(x'/\Delta) \bmod M$.

\begin{proposition}[Mod-$M$ distortion bound]
$\abs{x' - x} \leq M\Delta / 2$ for any mod-$M$ QIM embedding.
\end{proposition}

\begin{proof}
The maximum distance between $x$ and the nearest lattice point of class $s$
is bounded by the half-period of the sub-lattice spacing, which is $M\Delta/2$.
\end{proof}

\begin{remark}
StenoWav uses $M = 4$ for mask samples (2 bits per sample) and $M = 2$
for data samples (1 bit per sample).  The mask modulus $M = 4$ gives a
maximum per-sample distortion of $2\Delta$ versus $\Delta$ for $M = 2$.
\end{remark}

\subsection{Non-Negative QIM}
\label{sec:nonneg-qim}

In the frequency domain, embedding operates on magnitude coefficients
$\abs{X_k} \geq 0$.  Standard QIM can produce negative outputs when
$q^* < 0$, but negative magnitudes are physically meaningless and
would lose their sign through $\sqrt{\mathrm{Re}^2 + \mathrm{Im}^2}$
on re-extraction.

\begin{definition}[Non-negative mod-$M$ QIM]
Identical to Definition~2.4, except that $q_{\downarrow}$ is only
considered if $q_{\downarrow} \geq 0$.  If $q_{\downarrow} < 0$,
only $q_{\uparrow}$ is used.  Formally:
\[
    q^* = \begin{cases}
        \displaystyle\argmin_{q' \in \{q_\uparrow, q_\downarrow\}} \abs{q'\Delta - x}
            & \text{if } q_\downarrow \geq 0, \\[6pt]
        q_\uparrow & \text{otherwise.}
    \end{cases}
\]
\end{definition}

Extraction constrains the quantisation index similarly:
$\hat{s} = \max(0, \round(x'/\Delta)) \bmod M$.

\begin{proposition}
Non-negative mod-$M$ QIM preserves extraction correctness for all $x \geq 0$.
\end{proposition}

\begin{proof}
For $x \geq 0$, $q = \round(x/\Delta) \geq 0$, and the selected $q^* \geq 0$
by construction.  Since $q^* \equiv s \pmod{M}$ and
$\round(q^*\Delta / \Delta) = q^* \geq 0$, the extraction
$\max(0, q^*) \bmod M = s$.
\end{proof}


% ──────────────────────────────────────────────────────────────────
\section{Packet Architecture}
\label{sec:packet}
% ──────────────────────────────────────────────────────────────────

\subsection{Packet Structure}

Each packet embeds 2 bytes (16 bits) of payload using 9 non-contiguous
sample positions selected pseudorandomly:

\begin{center}
\begin{tabular}{@{}ccl@{}}
\toprule
Component & Count & Role \\
\midrule
Mask samples & 8 & Derive per-packet cipher via mod-4 QIM \\
Data region & 1 start index & 16 consecutive samples, 1 bit each via mod-2 QIM \\
\bottomrule
\end{tabular}
\end{center}

Total samples modified per packet: $8 + 16 = 24$.

\subsection{Index Selection}

Indices are selected via ChaCha20 CSPRNG, seeded from the master seed
$\sigma$ mixed with the packet number $p$:
\[
    \sigma_p = \sigma \xor (\underbrace{0, \ldots, 0}_{24}, p_0, p_1, \ldots, p_7),
\]
where $p_0, \ldots, p_7$ are the little-endian bytes of $p \in \N_0$.

From RNG stream $G(\sigma_p)$, 9 unique indices are drawn via rejection
sampling within the eligible pool (see Section~\ref{sec:rms}).  The first
8 indices are mask positions; the 9th is the data region start.

\subsection{Cipher Derivation}
\label{sec:cipher}

\begin{definition}[Cipher key]
A 16-bit value $K \in \{0, \ldots, 2^{16}-1\}$ derived from 8 mask samples.
Each mask sample $i$ contributes a 2-bit symbol $s_i \in \{0,1,2,3\}$
via mod-4 QIM extraction:
\[
    K = \sum_{i=0}^{7} s_i \cdot 4^{7-i}
      = \sum_{i=0}^{7} s_i \cdot 2^{2(7-i)}.
\]
Equivalently, $K$ is the concatenation of 8 two-bit values
$(s_0, s_1, \ldots, s_7)$ in big-endian order.
\end{definition}

\textbf{Encoding:}
A random cipher $K$ is generated from a \emph{separate} ChaCha20 stream
(seed byte 0 flipped: $\sigma' = \sigma \xor (0\text{xFF}, 0, \ldots, 0)$,
then mixed with $p$).  The symbols $s_i = (K \gg 2(7-i)) \mathbin{\&} 3$
are embedded into the mask samples via mod-4 QIM.

\textbf{Decoding:}
The decoder recomputes mask positions from $\sigma$ and $p$, extracts
mod-4 residues from each mask sample, and reconstructs $K$.

\subsection{Payload Embedding}

Given payload bytes $(b_\mathrm{hi}, b_\mathrm{lo})$ and cipher $K$:
\begin{enumerate}
    \item Concatenate into 16-bit word: $P = b_\mathrm{hi} \cdot 256 + b_\mathrm{lo}$.
    \item Encrypt: $P' = P \xor K$.
    \item Embed $P'$ bit-by-bit into the 16 consecutive samples starting
          at $\mathrm{data\_index}$, MSB first, using standard mod-2 QIM.
\end{enumerate}

\textbf{Extraction:}
\begin{enumerate}
    \item Extract 16 bits from the data region via mod-2 QIM.
    \item Reassemble into $P'$.
    \item Decrypt: $P = P' \xor K$ (XOR is self-inverse).
\end{enumerate}

\begin{proposition}[XOR cipher involution]
$\forall P, K: (P \xor K) \xor K = P$.
\end{proposition}

\subsection{Stream Framing}

An arbitrary byte stream is framed into packets:
\begin{itemize}
    \item Each packet carries 2 bytes.
    \item If the data length is odd, a zero byte is appended.
    \item Total packets: $N_\text{data} = \ceil{|\text{data}| / 2}$.
    \item The decoder uses the original length (from the header) to
          truncate any padding byte.
\end{itemize}


% ──────────────────────────────────────────────────────────────────
\section{Self-Contained Header}
\label{sec:header}
% ──────────────────────────────────────────────────────────────────

The first 7 packets embed a 14-byte metadata header, enabling decoding
with only the watermarked signal, seed, and configuration.

\begin{center}
\begin{tabular}{@{}ccl@{}}
\toprule
Bytes & Width & Content \\
\midrule
0--3 & 4 & Magic: \texttt{0x5354454E} (ASCII ``STEN'') \\
4    & 1 & Domain: $0 =$ Time, $1 =$ Frequency \\
5    & 1 & Reserved (0x00) \\
6--9 & 4 & \texttt{original\_len} (u32, big-endian) \\
10--13 & 4 & \texttt{num\_data\_packets} (u32, big-endian) \\
\bottomrule
\end{tabular}
\end{center}

\textbf{Packet layout:}
\[
    \underbrace{p_0, p_1, \ldots, p_6}_{\text{header (7 packets)}} \quad
    \underbrace{p_7, p_8, \ldots, p_{6+N}}_{\text{data ($N$ packets)}}
\]

The magic bytes serve as a validation check: if the wrong seed or delta
is used, the extracted magic will not match \texttt{STEN}, and decoding
fails immediately with a descriptive error.


% ──────────────────────────────────────────────────────────────────
\section{RMS-Adaptive Placement}
\label{sec:rms}
% ──────────────────────────────────────────────────────────────────

\subsection{Segment Analysis}

The signal is divided into non-overlapping segments of $L$ samples
(default $L = 4096$).  For each segment $j$:
\[
    \rms_j = \sqrt{\frac{1}{L_j} \sum_{n \in S_j} x_n^2},
\]
where $S_j$ is the set of sample indices in segment $j$ and $L_j = |S_j|$.

\subsection{Power Tier Classification}

Given thresholds $\theta_H > \theta_L > 0$ (defaults: $\theta_H = 0.10$,
$\theta_L = 0.01$):

\[
    \text{tier}(j) = \begin{cases}
        \textsc{High} & \text{if } \rms_j \geq \theta_H, \\
        \textsc{Mid}  & \text{if } \theta_L \leq \rms_j < \theta_H, \\
        \textsc{Skip} & \text{if } \rms_j < \theta_L.
    \end{cases}
\]

\subsection{Pool Construction}

From the classified segments, two pools are constructed:
\begin{itemize}
    \item \textbf{Mask pool}: All individual indices in High or Mid segments.
    \item \textbf{Data pool}: All indices $i$ in High or Mid segments
          such that $[i, i+16)$ lies entirely within that segment.
\end{itemize}

Packet index selection (Section~\ref{sec:packet}) draws from these pools,
ensuring no embedding occurs in quiet signal regions where QIM distortion
would be perceptible.

\begin{remark}[Threshold margin]
Thresholds should be set with margin above the QIM delta.  QIM modifies
samples by at most $\Delta$, shifting the segment RMS by approximately
$\Delta / \sqrt{L}$.  For $\Delta = 0.005$ and $L = 4096$:
$\Delta / \sqrt{L} \approx 7.8 \times 10^{-5}$,
which is negligible relative to $\theta_L = 0.01$.
\end{remark}


% ──────────────────────────────────────────────────────────────────
\section{Frequency-Domain Embedding}
\label{sec:freq}
% ──────────────────────────────────────────────────────────────────

\subsection{Analysis}

The signal is partitioned into $F$ non-overlapping frames of $W$ samples
(default $W = 1024$, must be a power of 2).  Incomplete tail frames are
discarded to ensure exact IDFT round-trip.

For each frame $f$:
\begin{align}
    X_f[k] &= \sum_{n=0}^{W-1} x_f[n] \, e^{-j 2\pi k n / W}, \quad k = 0, \ldots, W-1, \\
    |X_f[k]| &= \sqrt{\mathrm{Re}(X_f[k])^2 + \mathrm{Im}(X_f[k])^2}, \\
    \angle X_f[k] &= \mathrm{atan2}(\mathrm{Im}(X_f[k]), \mathrm{Re}(X_f[k])).
\end{align}

The DFT is computed via a radix-2 Cooley--Tukey FFT in~C.

\subsection{Magnitude Embedding}

Only bins in the range $[k_{\min}, k_{\max})$ are eligible (defaults:
$k_{\min} = 10$, $k_{\max} = 400$), excluding DC and very low frequencies
which carry perceptually critical content.

The eligible magnitudes from non-Skip frames are flattened into a single
array $\mathbf{m} = (m_0, m_1, \ldots)$.  Mask and data pools are
constructed over indices of $\mathbf{m}$ (analogous to Section~\ref{sec:rms}).

Mask samples use non-negative mod-4 QIM (Section~\ref{sec:nonneg-qim}).
Data samples use non-negative mod-2 QIM.

\subsection{Synthesis}

After embedding, the modified magnitudes $m'_i$ are scattered back to
their per-frame positions.  The time-domain signal is reconstructed:
\begin{enumerate}
    \item Rebuild the complex spectrum from modified magnitudes and
          \emph{original} phases:
          $X'_f[k] = m'_k \cdot e^{j \angle X_f[k]}$ for $k = 0, \ldots, W/2$.
    \item Enforce conjugate symmetry: $X'_f[W-k] = \overline{X'_f[k]}$
          for $k = 1, \ldots, W/2-1$.
    \item Apply IDFT: $x'_f[n] = \frac{1}{W} \sum_{k=0}^{W-1} X'_f[k] \, e^{j 2\pi kn / W}$.
\end{enumerate}

\begin{remark}
Phase preservation is critical: the decoder must compute the \emph{same}
DFT on the watermarked signal.  Since only magnitudes are modified, the
decoder's phase angles will differ slightly due to floating-point arithmetic,
but the magnitude round-trip is exact up to the DFT's numerical precision.
\end{remark}

\subsection{Frequency-Domain QIM Delta}

Frequency magnitudes can be much larger than normalised PCM samples.
A separate delta $\Delta_f$ (default 0.05) controls the frequency-domain
step size, independent of the time-domain $\Delta_t$ (default 0.005).

The SNR impact in the frequency domain is analogous:
\[
    \snr_f = 20 \log_{10}\!\left(\frac{\rms(\mathbf{m})}{\Delta_f / 2}\right) \text{ dB}.
\]


% ──────────────────────────────────────────────────────────────────
\section{Dither}
\label{sec:dither}
% ──────────────────────────────────────────────────────────────────

\subsection{Motivation}

QIM places modified samples on a regular lattice.  An attacker performing
statistical analysis could detect that some samples sit on a grid while
neighbouring samples do not.  Dither counters this by adding small
perturbations to \emph{all} non-payload samples, making the entire
signal appear uniformly noisy.

\subsection{Time-Domain Dither}

For each sample index $n$:
\begin{enumerate}
    \item Generate $u_n \sim \mathcal{U}(-\alpha, +\alpha)$ from a ChaCha20 CSPRNG
          seeded with $\sigma_\text{dither} = \sigma \xor (0, 0, \text{0xFF}, 0, \ldots)$
          (byte 2 flipped for stream independence).
    \item If $n \notin \mathcal{P}$ (the set of protected indices):
          $x_n \leftarrow \mathrm{clamp}(x_n + u_n, -1, 1)$.
    \item If $n \in \mathcal{P}$: leave $x_n$ unchanged (but still
          consume $u_n$ to maintain RNG synchronisation).
\end{enumerate}

Default intensity: $\alpha = 2 / 32768 \approx 6.1 \times 10^{-5}$
(two least-significant bits of 16-bit PCM).

The protected set $\mathcal{P}$ is the union of all mask indices and
data regions $[\mathrm{data\_index}, \mathrm{data\_index} + 16)$ across
all packets.

\subsection{Frequency-Domain Dither}

For frequency-domain embedding, dither is applied to magnitude bins
\emph{before} IDFT synthesis:
\begin{enumerate}
    \item Generate $u_i \sim \mathcal{U}(-\alpha_f, +\alpha_f)$ per flat-magnitude index.
    \item If $i \notin \mathcal{P}_f$: $m_i \leftarrow \max(0, m_i + u_i)$.
\end{enumerate}

The intensity is auto-scaled: $\alpha_f = \alpha \cdot (\Delta_f / \Delta_t)$.

\begin{proposition}[Dither does not affect extraction]
Since the decoder reads only from indices in $\mathcal{P}$ (or $\mathcal{P}_f$),
and dither skips all protected indices, the embedded data is unaffected.
\end{proposition}

\subsection{Audibility Analysis}

The dither noise power per sample is:
\[
    \sigma_\text{dither}^2 = \frac{\alpha^2}{3}
    \quad\Rightarrow\quad
    \sigma_\text{dither} = \frac{\alpha}{\sqrt{3}}.
\]

For the default $\alpha = 6.1 \times 10^{-5}$:
$\sigma_\text{dither} \approx 3.5 \times 10^{-5}$.

For a typical signal with $\rms(x) = 0.2$:
\[
    \snr_\text{dither} = 20 \log_{10}\!\left(\frac{0.2}{3.5 \times 10^{-5}}\right) \approx 75 \text{ dB}.
\]

This is well above the 50\,dB empirical threshold measured on real audio
(My Own Summer, 16-bit 44.1kHz) and far above the 20\,dB threshold
generally considered transparent for audio watermarking.


% ──────────────────────────────────────────────────────────────────
\section{Quality Metrics}
\label{sec:quality}
% ──────────────────────────────────────────────────────────────────

Given original signal $\mathbf{x}$ and watermarked signal $\mathbf{x}'$,
both of length $N$:

\begin{align}
    \mathbf{d} &= \mathbf{x}' - \mathbf{x} && \text{(difference signal)} \\[4pt]
    \snr &= 20 \log_{10}\!\left(\frac{\rms(\mathbf{x})}{\rms(\mathbf{d})}\right) \text{ dB}
    && \text{where } \rms(\mathbf{v}) = \sqrt{\frac{1}{N}\sum_{n} v_n^2} \\[4pt]
    \delta_{\max} &= \max_n \abs{d_n} && \text{(worst-case distortion)} \\[4pt]
    \bar{\delta} &= \frac{1}{N}\sum_n \abs{d_n} && \text{(mean distortion)} \\[4pt]
    S_\text{diff} &= \max_k \abs{|\F(\mathbf{x}')[k]| - |\F(\mathbf{x})[k]|}
    && \text{(max spectral magnitude difference)}
\end{align}

Empirical measurements on real audio (My Own Summer, 44.1kHz stereo):
\begin{center}
\begin{tabular}{@{}lccc@{}}
\toprule
Scenario & SNR (dB) & $\delta_{\max}$ & $S_{\text{diff}}$ \\
\midrule
Embedding only ($\Delta_t = 0.005$) & $> 30$ & $< 0.01$ & $< 5.0$ \\
Embedding + dither ($\alpha = 2\,\text{LSB}$) & $> 30$ & $< 0.01$ & $< 5.0$ \\
Pure dither ($\alpha = 2\,\text{LSB}$) & $> 50$ & $< 10^{-4}$ & $< 0.5$ \\
\bottomrule
\end{tabular}
\end{center}


% ──────────────────────────────────────────────────────────────────
\section{Security Analysis}
\label{sec:security}
% ──────────────────────────────────────────────────────────────────

\subsection{Key Space}

The master seed $\sigma$ is 32 bytes (256 bits), giving a key space of
$2^{256}$.  The CSPRNG (ChaCha20) is a well-studied stream cipher with
no known practical attacks below brute force.

\subsection{Extraction Requirements}

An attacker attempting to extract the hidden data needs:
\begin{enumerate}
    \item \textbf{The seed $\sigma$}: to locate mask and data sample positions.
          Without $\sigma$, the attacker cannot determine which of the millions
          of samples carry information.
    \item \textbf{The QIM delta $\Delta$}: to extract mod-4 residues from mask
          samples and mod-2 bits from data samples.  The wrong $\Delta$ yields
          incorrect quantisation indices and thus garbage output.
    \item \textbf{Configuration parameters}: domain, RMS thresholds, segment
          size, frequency bin range.  These determine the eligible pools and
          thus which sample positions the CSPRNG selects from.
\end{enumerate}

\subsection{Statistical Resistance}

\begin{itemize}
    \item \textbf{Without dither}: Modified samples lie on a QIM lattice
          ($\Delta$-spaced grid), while unmodified samples retain their
          original distribution.  A statistical test comparing sample-value
          distributions could detect the lattice pattern.
    \item \textbf{With dither}: All non-payload samples receive random
          perturbations, obscuring the boundary between modified and
          unmodified regions.  The entire signal exhibits uniform micro-noise.
\end{itemize}

\subsection{Per-Packet Cipher}

Each packet generates an independent cipher from a separate ChaCha20 stream
(byte 0 of seed flipped).  This means:
\begin{itemize}
    \item Identical payloads in different packets produce different embedded
          bit patterns.
    \item Knowledge of one packet's cipher does not reveal others.
    \item The header magic provides a quick wrong-key check (probability of
          false positive: $\approx 2^{-32}$).
\end{itemize}


% ──────────────────────────────────────────────────────────────────
\section{WAV I/O}
\label{sec:wav}
% ──────────────────────────────────────────────────────────────────

\subsection{Sample Representation}

Internal: $x \in [-1.0, 1.0]$ as \texttt{f64}.

\textbf{Decode (file $\to$ internal):}
\[
    x = \frac{\text{i16\_sample}}{32768.0}.
\]

\textbf{Encode (internal $\to$ file):}
\[
    \text{i16\_sample} = \mathrm{clamp}\!\left(\round(x \cdot 32768.0),\; -32768,\; 32767\right).
\]

\begin{remark}[Symmetry constant]
Both directions use $32768.0$.  An earlier design used $32767.0$ for encoding,
causing asymmetric quantisation error ($\sim 4 \times 10^{-5}$ vs
$\sim 3 \times 10^{-5}$).  Using the same constant ensures the round-trip
error is symmetric and bounded by $1/(2 \cdot 32768) \approx 1.5 \times 10^{-5}$.
\end{remark}

\begin{theorem}[QIM survives i16 round-trip]
\label{thm:wav-roundtrip}
For $\Delta \gg 1/32768 \approx 3.05 \times 10^{-5}$, QIM-embedded bits
survive the \texttt{f64} $\to$ \texttt{i16} $\to$ \texttt{f64} round-trip.
\end{theorem}

\begin{proof}[Proof sketch]
The i16 round-trip introduces at most $\epsilon = 1/(2 \cdot 32768)$
error per sample.  For extraction to succeed, the watermarked sample
$x'$ must remain in the Voronoi region of its QIM lattice point $q^*\Delta$.
The Voronoi radius is $\Delta/2$.  If $\epsilon < \Delta/2$, i.e.
$\Delta > 2\epsilon \approx 3.05 \times 10^{-5}$, extraction is correct.
At $\Delta = 0.005$, the margin is $> 160\times$.
\end{proof}


% ──────────────────────────────────────────────────────────────────
\section{FFT Implementation}
\label{sec:fft}
% ──────────────────────────────────────────────────────────────────

The DFT/IDFT is implemented in C using the radix-2 Cooley--Tukey algorithm:

\begin{itemize}
    \item In-place computation with bit-reversal permutation.
    \item Input length must be a power of 2.
    \item Split real/imaginary arrays (\texttt{double *re, *im}).
    \item IDFT via: conjugate $\to$ FFT $\to$ conjugate $\to$ scale by $1/N$.
    \item Complexity: $O(N \log N)$.
\end{itemize}

Additional C-level utilities: RMS computation, complex magnitude
($\sqrt{\mathrm{Re}^2 + \mathrm{Im}^2}$), complex phase
($\mathrm{atan2}(\mathrm{Im}, \mathrm{Re})$).

Rust FFI bindings validate all preconditions (power-of-2 length,
matching array lengths) before calling C, returning typed errors
via \texttt{FftError}.


% ──────────────────────────────────────────────────────────────────
\section{Complete Pipeline}
\label{sec:pipeline}
% ──────────────────────────────────────────────────────────────────

\subsection{Encoding}

\begin{algorithm}[H]
\caption{StenoWav Encoding Pipeline}
\begin{algorithmic}[1]
\Require WAV file path, data bytes $D$, seed $\sigma$, config $C$
\Ensure Watermarked WAV file
\State $\mathbf{x} \gets \text{read\_wav}(\text{path})[\text{channel}]$
\State $\text{segments} \gets \text{analyze\_rms}(\mathbf{x}, C)$
\State $(\text{mask\_pool}, \text{data\_pool}) \gets \text{build\_pools}(\text{segments})$
\State $N_D \gets \lceil |D| / 2 \rceil$ \Comment{data packets}
\State $H \gets \text{serialize\_header}(\text{domain}, |D|, N_D)$ \Comment{14 bytes}
\State $\text{payload} \gets H \| \text{pad}(D)$
\For{$p = 0$ \textbf{to} $6 + N_D$}
    \State $\text{idx} \gets \text{select\_indices}(\sigma, p, \text{pools})$
    \State $K_p \gets \text{generate\_cipher}(\sigma, p)$
    \State Embed $K_p$ into mask samples via mod-4 QIM
    \State Embed $\text{payload}[2p:2p{+}2] \xor K_p$ into data region via mod-2 QIM
\EndFor
\If{dither enabled}
    \State $\mathcal{P} \gets \text{collect\_protected\_indices}()$
    \State Apply dither to $\mathbf{x}[n]$ for $n \notin \mathcal{P}$
\EndIf
\State $\text{write\_wav}(\text{output}, \mathbf{x})$
\end{algorithmic}
\end{algorithm}

\subsection{Decoding}

\begin{algorithm}[H]
\caption{StenoWav Decoding Pipeline}
\begin{algorithmic}[1]
\Require Watermarked WAV path, seed $\sigma$, config $C$
\Ensure Extracted data bytes or error
\State $\mathbf{x}' \gets \text{read\_wav}(\text{path})[\text{channel}]$
\State Recompute segments and pools (same as encoding)
\For{$p = 0$ \textbf{to} $6$} \Comment{Extract header}
    \State Extract cipher $K_p$ from mask samples
    \State Extract payload $\xor$ $K_p$ from data region
\EndFor
\State Validate magic = \texttt{STEN}; if wrong, return error
\State Parse $|D|$ and $N_D$ from header
\For{$p = 7$ \textbf{to} $6 + N_D$} \Comment{Extract data}
    \State Extract cipher and payload (same as header)
\EndFor
\State Truncate to $|D|$ bytes; return
\end{algorithmic}
\end{algorithm}


% ──────────────────────────────────────────────────────────────────
\section{API Reference}
\label{sec:api}
% ──────────────────────────────────────────────────────────────────

\subsection{Primary Interface}

\begin{lstlisting}[language={}]
Steno::encode(input_path, output_path, data, key, config) -> Result<(), StenoError>
Steno::decode(watermarked_path, key, config)              -> Result<Vec<u8>, StenoError>
\end{lstlisting}

\subsection{In-Memory Interface}

\begin{lstlisting}[language={}]
Steno::encode_samples(samples, data, key, config) -> Result<(), StenoError>
Steno::decode_samples(samples, key, config)        -> Result<Vec<u8>, StenoError>
\end{lstlisting}

\subsection{Configuration Parameters}

\begin{center}
\begin{tabular}{@{}llcl@{}}
\toprule
Parameter & Type & Default & Description \\
\midrule
\texttt{domain} & enum & Time & Time or Frequency embedding \\
\texttt{time\_delta} & f64 & 0.005 & QIM $\Delta$ for PCM samples \\
\texttt{freq\_delta} & f64 & 0.05 & QIM $\Delta$ for magnitudes \\
\texttt{rms\_high} & f64 & 0.10 & High power tier threshold \\
\texttt{rms\_low} & f64 & 0.01 & Skip tier threshold \\
\texttt{rms\_segment\_size} & usize & 4096 & Samples per RMS segment \\
\texttt{freq\_window\_size} & usize & 1024 & FFT window (power of 2) \\
\texttt{freq\_min\_bin} & usize & 10 & Lowest eligible frequency bin \\
\texttt{freq\_max\_bin} & usize & 400 & Highest eligible bin (excl.) \\
\texttt{dither\_enabled} & bool & false & Enable noise dither \\
\texttt{dither\_intensity} & f64 & $\sim$6.1e-5 & Max perturbation ($\sim$2 LSB) \\
\texttt{channel} & usize & 0 & Audio channel (0 = left) \\
\bottomrule
\end{tabular}
\end{center}


% ──────────────────────────────────────────────────────────────────
\section{Capacity Analysis}
\label{sec:capacity}
% ──────────────────────────────────────────────────────────────────

Each packet carries 2 bytes and uses $\sim$24 sample positions.  For a
message of $B$ bytes:
\begin{align}
    N_\text{data} &= \ceil{B / 2}, \\
    N_\text{total} &= 7 + N_\text{data}, \\
    \text{samples used} &= 24 \cdot N_\text{total}.
\end{align}

To avoid packet collisions (overlapping random indices), the eligible pool
should be $\gg 24 \cdot N_\text{total}$.  A conservative guideline:
\[
    |\text{pool}| \geq 100 \times 24 \cdot N_\text{total}.
\]

For typical 44.1kHz stereo audio (3 minutes $\approx$ 7.9M samples per channel),
and assuming 50\% of segments are eligible, the pool size is $\sim$4M,
supporting $\sim$6900 packets $\approx$ 13.8\,KB of hidden data.


% ──────────────────────────────────────────────────────────────────
\section{Future Directions}
\label{sec:future}
% ──────────────────────────────────────────────────────────────────

\begin{enumerate}
    \item \textbf{Variable packet sizes}: Expand beyond 9 samples per
          packet (e.g., 17 for 32-bit cipher + 32 data bits).
    \item \textbf{Higher bit density}: Use higher-order QIM ($M > 4$
          for masks, $M > 2$ for data) to increase bits per sample.
    \item \textbf{AWGN robustness optimisation}: Sweep $\Delta$, RMS
          thresholds, and bit density across thousands of noise samples
          to find the Pareto frontier of robustness vs.\ imperceptibility.
    \item \textbf{Compression robustness}: Test and harden against MP3,
          AAC, and Opus transcoding.
    \item \textbf{Streaming / real-time}: Adapt the frame-based
          architecture for low-latency applications.
    \item \textbf{Error correction}: Add forward error correction (e.g.,
          Reed--Solomon) for graceful degradation under channel noise.
    \item \textbf{Multiple channel embedding}: Spread data across stereo
          channels for increased capacity and redundancy.
\end{enumerate}


% ──────────────────────────────────────────────────────────────────
\appendix
\section{Module Map}
\label{app:modules}
% ──────────────────────────────────────────────────────────────────

\begin{center}
\begin{tabular}{@{}ll@{}}
\toprule
File & Responsibility \\
\midrule
\texttt{src/steno.rs} & Public API: \texttt{Steno}, \texttt{StenoConfig}, \texttt{StenoError} \\
\texttt{src/embed.rs} & Dual-domain engine, RMS analysis, header, \texttt{encode\_full}/\texttt{decode\_full} \\
\texttt{src/packet.rs} & Packet structure, index selection, encode/decode per packet \\
\texttt{src/cipher.rs} & Cipher key type, generation, mod-4 QIM embed/extract \\
\texttt{src/qim.rs} & Core QIM primitives (mod-2) \\
\texttt{src/dither.rs} & Dither engine, protected index collection, quality metrics \\
\texttt{src/fft.rs} & Rust FFI bindings to C FFT library \\
\texttt{src/wav.rs} & 16-bit PCM WAV reader/writer \\
\texttt{csrc/math\_ops.c} & C math library: FFT, IFFT, RMS, complex ops \\
\texttt{build.rs} & Compiles C via \texttt{cc} crate \\
\bottomrule
\end{tabular}
\end{center}

\end{document}
